\section{Solver集合における離散構造と連続構造の対応}

ECAでは、公理体系を単なる公理の集合として扱うのではなく、それ自体を数学的対象として扱うことから出発する。すなわち、

\[
\text{命題}
\longrightarrow
\text{公理体系}
\longrightarrow
\text{公理体系の集合}
\longrightarrow
\text{公理体系間の構造}
\]

という対象化を行う。

この対象化によって得られる公理体系の集合を基礎として、複数の公理体系およびそれらの間に存在する構造をまとめて扱うための対象としてSolver集合 $\mathcal{Z}$ を考える。

ここでいうSolver集合は、特定の選択主体や評価主体を意味するものではない。Solver集合は、公理体系を対象化することによって得られる数学的構造を記述するための対象である。

\subsection{Solverの合成とSolver集合}

複数の公理体系を個別に扱うだけでは、それらの間に存在する共通部分、差異、包含関係、あるいはそれらを組み合わせた構造を一般的に記述することが難しい。

そこで、複数の公理体系を一つの対象として扱う構造を導入する。

公理体系を $\Sigma_1,\Sigma_2,\ldots,\Sigma_n$ とすると、それらから構成される集合を

\[
\mathcal{Z}
=
\{\Sigma_1,\Sigma_2,\ldots,\Sigma_n\}
\]

と表すことができる。

さらに、複数の公理体系に対応する構造を同時に扱う操作を考えることで、それらの合成を一つの数学的対象として記述することができる。

ここでいうSolverの合成は、Solverを独立した主体として組み合わせることを意味するのではなく、複数の公理体系に対応する構造を一つの対象として扱うことを意味する。

したがって、Solverの合成はSolver集合の構成と分離された概念ではなく、

\[
\text{個々の公理体系}
\longrightarrow
\text{公理体系の集合}
\longrightarrow
\text{合成された構造}
\]

という関係の中で理解される。

この構造によって、個々の公理体系だけでなく、公理体系の集合そのもの、およびその集合上に定義される構造を対象とすることが可能になる。

\subsection{離散構造と連続構造}

Solver集合 $\mathcal{Z}$ を扱う際には、その構造を離散的に記述する場合と連続的に記述する場合を区別することができる。

離散的な記述では、個々の公理体系や公理選択を要素として扱い、それらの包含関係、共通部分、差異、組合せなどを集合論的に記述する。

一方、連続的な記述では、公理体系そのものを連続的な対象とみなすのではなく、公理体系またはSolver集合上に定義された量、写像、測度、パラメータなどを用いて、その構造を連続的な形式で表現する。

したがって、ここでいう離散構造と連続構造は、単純に「有限集合」と「実数」の対比を意味するものではない。

むしろ、

\[
\text{Solver集合の構造}
\]

に対して、

\[
\text{離散的な構造表現}
\]

と

\[
\text{連続的な構造表現}
\]

という異なる表現形式を考えることができる、という関係にある。

この区別は、ECAにおける離散式と連続式の区別とも関係する。ただし、離散式・連続式という語は二つの異なる文脈で現れるため、それらを同一視してはならない。

一つは、Solver集合そのものにおける離散構造と連続構造である。

もう一つは、ある構造に対して評価を行う際の評価形式としての離散式と連続式である。

したがって、

\[
\boxed{
\text{構造としての離散性・連続性}
}
\]

と

\[
\boxed{
\text{評価形式としての離散式・連続式}
}
\]

は区別して扱う必要がある。

\subsection{離散構造から連続構造への正規化}

離散的に記述されたSolver構造と、連続的に記述されたSolver構造との間に対応を与えるため、正規化写像を考える。

一般に、離散的な構造を $D$、連続的な構造を $C$ とすると、

\[
N:D\longrightarrow C
\]

という写像によって、離散的な構造を連続的な表現へ写すことができる。

ただし、現在のECAにおいて重要なのは、正規化そのものをECAの中心原理とすることではない。

正規化は、Solver集合における同一の構造を異なる形式で記述するための対応付けとして位置づけられる。

すなわち、Solver集合 $\mathcal{Z}$ に対して、その構造を離散的に記述する形式と連続的に記述する形式との間に、

\[
N:
\mathcal{Z}_{\mathrm{disc}}
\longrightarrow
\mathcal{Z}_{\mathrm{cont}}
\]

という対応を考えることができる。

ここで $\mathcal{Z}_{\mathrm{disc}}$ と $\mathcal{Z}_{\mathrm{cont}}$ は、必ずしもSolver集合そのものを二つに分割することを意味しない。

これらは、同一のSolver集合 $\mathcal{Z}$ に対して与えられる異なる構造表現として扱うことができる。

したがって、離散的な構造と連続的な構造は、互いに独立した二つの理論としてではなく、同一の数学的対象に対する異なる形式として扱うことができる。

\subsection{離散構造と連続構造の弱可換性}

離散的な構造上で定義される変換または操作を $F$ とし、連続的な構造上で対応する変換または操作を $G$ とする。

このとき、

\[
\begin{array}{ccc}
D & \xrightarrow{F} & D \\
\downarrow N & & \downarrow N \\
C & \xrightarrow{G} & C
\end{array}
\]

という図式を考えることができる。

この図式において、厳密な可換性

\[
N\circ F
=
G\circ N
\]

が成立するとは限らない場合でも、両者が適切な意味で対応するのであれば、

\[
N\circ F
\approx
G\circ N
\]

という弱い可換性を考えることができる。

ここで $\approx$ が具体的に何を意味するかは、対象とするSolver構造および正規化写像 $N$ に応じて定義されなければならない。

したがって、弱可換図式そのものを一般的なECAの定理として直ちに主張するのではなく、離散的構造と連続的構造の間に成立し得る対応を記述するための形式として位置づける。

Solver集合 $\mathcal{Z}$ に対してこの対応を記述する場合には、

\[
\begin{array}{ccc}
\mathcal{Z}_{\mathrm{disc}}
& \xrightarrow{F} &
\mathcal{Z}_{\mathrm{disc}}
\\
\downarrow N
&&
\downarrow N
\\
\mathcal{Z}_{\mathrm{cont}}
& \xrightarrow{G} &
\mathcal{Z}_{\mathrm{cont}}
\end{array}
\]

と表すことができる。

このとき、

\[
N\circ F
\approx
G\circ N
\]

は、Solver集合に対する離散的な操作を連続的な表現へ移した場合と、最初から連続的な表現上で対応する操作を行った場合との間に、どの程度の対応が存在するかを記述する。

\subsection{評価式としての離散式・連続式との区別}

以上で扱った離散構造と連続構造は、ECAにおける評価式としての離散式・連続式とは区別される。

ECAでは、構造に対する評価量を離散的に記述する場合、

\[
T_{\mathrm{overlay}}=\frac{1}{N}\sum_{i=1}^{N}B_iI_iK_i
\]

と表すことができ、連続的に記述する場合には、

\[
T_{\mathrm{overlay}}=\int_0^1 B(x)I(x)K(x)\,dx
\]

と表すことができる。

ここで離散式と連続式は、Solver集合そのものの構造を定義するものではなく、構造に対して評価量を与えるための形式として扱われる。

したがって、ECAでは次の二つを明確に区別する。

\[
\text{Solver集合における離散構造・連続構造}
\]

と

\[
\text{構造に対する評価式としての離散式・連続式}
\]

前者は公理体系を対象化した数学的構造そのものに関する問題であり、後者はその構造をどのような形式で評価するかに関する問題である。

この区別によって、離散性と連続性をECAの本質そのものとして扱うのではなく、公理体系を対象化して得られた構造を記述・評価するための異なる形式として整理することができる。

\subsection{位置づけ}

以上の構造は、旧来のECAにおいて扱われていたSolverの合成、離散構造と連続構造の対応、離散から連続への正規化、および弱可換図式を、現在のECAの枠組みに合わせて再構成するための基礎となる。

現在のECAでは、公理体系の対象化が先にあり、その対象化からSolver集合が構成される。

したがって、

\[
\text{公理体系の対象化}
\longrightarrow
\mathcal{Z}
\longrightarrow
\begin{cases}
\text{離散構造}\\
\text{連続構造}
\end{cases}
\]

という順序が基本となる。

その上で、離散構造と連続構造の対応を正規化写像によって記述し、必要に応じて弱可換図式によってその対応関係を表現する。

この構造は、ECAにおける評価式そのものを置き換えるものではない。むしろ、Solver集合という対象をどのような構造表現によって扱うことができるかを明確化するための数学的枠組みとして位置づけられる。
